\newpage
\section{Analyse des Ausgangszustands und Ableitung des Kriterienkatalogs}
\label{sec:bedrohungsmodell}

Der Maßstab für das zu konstruierende Artefakt stammt aus zwei Quellen, die beide in der
Fallstudie liegen. Die erste ist das Konzept, mit dem die Scripting-Engine heute betrieben wird,
und die Vertrauensannahmen, die daraus folgen. Die zweite ist der produktive Skriptbestand, der
festlegt, welche Schnittstelle ein neues Modell mindestens bereitstellen muss. Aus beidem ergeben
sich die relevanten Angriffswege, und aus diesen der Kriterienkatalog, an dem sich der Entwurf in
\autoref{sec:Entwurf} und die Evaluation in \autoref{sec:Evaluation} messen lassen.

\subsection{Ausgangszustand und Vertrauensannahmen}
\label{sec:systemkontext}

Die Scripting-Engine der Fallstudie ist ein mandantenfähiger Execution Service. Sie führt
JavaScript auf Node.js aus und verarbeitet damit Webhooks, die Ereignisse im \ac{DMS} auslösen
\autocite{dvelop_ag_dvelop_2026}. Den Code hinterlegt der Mandant selbst. Er wird versioniert
gespeichert und über viele Ausführungen hinweg unverändert wiederverwendet. Künftig soll
dieselbe Engine zusätzlich als Werkzeug eines KI-Agenten dienen, der kurze Skripte zur Laufzeit
erzeugt und ausführen lässt \autocite{wang_executable_2024}. Diese
Fähigkeit besteht im Produkt noch nicht. Sie gehört deshalb nicht zum Ausgangszustand, sondern
ist eine Anforderung an das Artefakt, das beide Profile tragen muss. Betriebsdaten liegen für das
Agentenprofil nicht vor, weshalb alle Aussagen dazu aus dem Interaktionsmuster eines Agenten
hergeleitet und als Annahme gekennzeichnet sind.

Zwei Ressourcen teilen sich die Ausführungen heute, und an beiden muss eine Mandantengrenze
verlaufen. Die erste ist der Host-Kernel. Der Execution Service schaltet seinen Runner-Pods
dafür eine gVisor-Runtime-Klasse vor und ergänzt sie um einen gehärteten Sicherheitskontext, bei
dem Privilege Escalation unterbunden ist, sämtliche Capabilities entzogen sind, das
seccomp-Standardprofil greift und kein Service-Account-Token eingebunden wird.\footnote{Die
Angaben zur Konfiguration des Execution Service beruhen auf einer eigenen Auswertung des
Quellcodes der Plattform, Stand September 2026.} Der rohe Zugriff auf die Systemaufrufe des
Hosts, den generische Containerisierung offenlegt \autocite{souppaya_application_2017}, wird
dadurch verkleinert, aber nicht beseitigt.

Welches Skript eine Instanz ausführt, entscheidet sich erst beim Aufruf. Das Skript wird je
Ausführung über eine vorsignierte URL aus einem Objektspeicher nachgeladen, weshalb eine Instanz
nicht an ein Skript gebunden ist, sondern nacheinander beliebige Skripte ihres Mandanten
ausführt. Sein Eintrag im Modul-Zwischenspeicher wird dabei vor jeder Ausführung verworfen, sodass
auch das Übersetzen bei jeder Ausführung erneut anfällt.

Die zweite geteilte Ressource liegt innerhalb dieser Sandbox. Jeder Mandant erhält einen eigenen
Execution Service, sodass Instanzen mandantenübergreifend nicht wiederverwendet werden. Der
Skriptcode wird unmittelbar in den laufenden Node.js-Prozess geladen, ohne einen eigenen
Unterprozess je Ausführung, und es läuft höchstens eine Ausführung zur Zeit. Aufeinanderfolgende
Ausführungen desselben Mandanten teilen sich damit den Prozessspeicher. Das temporäre Verzeichnis
wird vor jeder Ausführung zurückgesetzt, Heap und globale Variablen bleiben bestehen. Ein
Node.js-Prozess je Mandant statt je Ausführung ist der Ausgangspunkt, gegen den sich jedes
geprüfte Modell messen lassen muss.

Aus diesem Konzept folgen die Schutzziele. \emph{Vertraulichkeit} und \emph{Integrität}
verlangen, dass eine Ausführung Daten, Code und Verarbeitung anderer Mandanten weder einsehen
noch verändern kann. Die \emph{Verfügbarkeit} verlangt, dass sie die Bedienung anderer Mandanten
nicht beeinträchtigt. Hinzu tritt die \emph{Zustandsfreiheit}, die verlangt, dass eine Ausführung
keinen Zustand an eine spätere weitergibt.

Das vierte Ziel verdient eine genauere Betrachtung, weil sein Charakter am Lastprofil hängt. In
der Webhook-Verarbeitung stammen alle Skripte eines Mandanten von ihm selbst und bilden eine
gemeinsame Vertrauensdomäne. Übernommener Zustand stört dort die folgenden Ausführungen,
überschreitet aber keine Vertrauensgrenze, weshalb Zustandsfreiheit hier eine Anforderung an die
Zuverlässigkeit ist. Im Agentenprofil gilt das nicht mehr. Zwei aufeinanderfolgende Ausführungen
desselben Mandanten entstehen aus verschiedenen Kontexten und können nach V2 von verschiedenen
Dritten beeinflusst sein, sind einander also nicht wechselseitig vertrauenswürdig. Der
Zustandsübertrag wird damit auch innerhalb eines Mandanten zur Vertraulichkeitsverletzung. Die
Vertrauenseinheit verschiebt sich vom Mandanten auf die einzelne Ausführung, und die
Isolationsgrenze muss dieser Verschiebung folgen. Diese Verschiebung ist der Kern des Problems,
das die vorliegende Arbeit löst.

Damit lassen sich die Vertrauensannahmen festhalten.

\begin{description}
  \item[V1: Der ausgeführte Code ist nicht vertrauenswürdig.] Im Webhook-Profil hinterlegt der
    Mandant den Code selbst, kennt ihn und verantwortet ihn. Aus Sicht der Plattform bleibt er
    fremder Code, weil ein Mandant nach V4 bösartig sein kann. Im Agentenprofil entfällt
    zusätzlich die Annahme, der Mandant kenne den Code, denn dieser entsteht erst zur Laufzeit
    und wird vor seinem einzigen Lauf von keinem Menschen gelesen
    \autocite{wang_executable_2024}. Er gilt in beiden Fällen nicht als fehlerhaft, sondern als
    potenziell aktiv bösartig.

  \item[V2: Im Agentenprofil ist der Code über seine Eingaben von Dritten beeinflussbar.] Das
    Modell erzeugt den Code aus einem Kontext, der Inhalte anderer Personen enthält, und solche
    Inhalte lassen sich gezielt so gestalten, dass sie das Verhalten des Modells steuern
    \autocite{greshake_not_2023}. Der Urheber des schädlichen Codes muss deshalb nicht der
    Mandant sein. Dass von Agenten erzeugter Code risikobehaftete Operationen enthalten kann, ist
    empirisch belegt \autocite{guo_redcode_2024}.

  \item[V3: Die Plattform ist vertrauenswürdig.] Host-Kernel, gVisor, Container-Orchestrierung,
    Laufzeitbinärdatei und Execution Service gelten als korrekt und nicht kompromittiert.
    Angriffe, die diese Bestandteile auf anderem Weg unterwandern, etwa über die
    Kubernetes-Steuerungsebene, liegen außerhalb der Betrachtung. Das codeerzeugende Modell
    gehört nach V5 nicht dazu.

  \item[V4: Die Mandanten misstrauen einander.] Es besteht keine Annahme, dass ein Mandant die
    Ausführungen eines anderen nicht angreift.

  \item[V5: Das codeerzeugende Modell liegt außerhalb der Kontrolle der Plattform.] Der Mandant
    kann ein eigenes Modell anbinden und kontrolliert dann Gewichte, System prompt und
    vorgelagerte Prüfungen vollständig. Aussagen über das Alignment dieses Modells, über seine
    Widerstandsfähigkeit gegen die in V2 beschriebene Beeinflussung oder über die Einhaltung
    einer vorgegebenen Schnittstelle sind deshalb nicht möglich.
\end{description}

V5 hat eine unmittelbare Folge für den Entwurf. Eine Beschränkung der Fähigkeiten, die einem
Skript zur Verfügung stehen, kann nicht über System prompt oder Tool Schema erfolgen, weil beides
dem Mandanten gehört. Sie muss von der Ausführungsumgebung durchgesetzt werden. Für das
Agentenprofil verbleibt der Plattform damit keine Schutzmaßnahme außerhalb der Isolationsgrenze.
Die jüngere Forschung zur Sicherheit agentischer Systeme betrachtet diese Frage bislang
überwiegend aus der Perspektive des Modells und weniger aus der der Systemisolation
\autocite{christodorescu_systems_2025}.

\subsection{Struktur des produktiven Skriptbestands}
\label{sec:bestandsanalyse}

Ein Ausführungsmodell, das den heutigen Skripten die benötigte Schnittstelle entzieht, ist
unabhängig von seinen Isolationseigenschaften nicht einführbar. Dieser Abschnitt erhebt deshalb,
welchen Teil der Node.js-Schnittstelle der Bestand tatsächlich braucht.

\subsubsection{Datengrundlage und Vorgehen}

Ausgewertet wurden 1.718 Mandanten mit 14.575 Skriptversionen, die sich auf 8.310 inhaltlich
verschiedene Skripte und 4.768 Importe verteilen. Vollständig minifizierte Inhalte sind
ausgeschlossen, weil sich ihre Importe statisch nicht zuverlässig bestimmen lassen. Der Filter
entfernt 269 Inhalte, 519 Skriptversionen und 25 Mandanten, die ausschließlich solche Skripte
vorhalten. Eine Gegenrechnung ohne den Filter verschiebt die Kennzahlen um gut einen
Prozentpunkt und ändert die Aussage nicht.

Jedes Skript wird anhand seiner Importe klassifiziert, wobei die härteste Einschränkung zählt.
Ein Mandant erhält die Einstufung seines härtesten Skripts und nicht den Durchschnitt. Die
Analyse ist statisch und sieht daher weder dynamisch zusammengesetzte Modulnamen noch gebündeltes
Kompilat. Von den Skripten fallen 24 als Parse-Fehler an und werden vorsorglich als blockierend
geführt.

Ein Import allein sagt nichts darüber, was ein npm-Paket seinerseits benötigt. Für die Pakete mit
nennenswertem Mandantengewicht wurde deshalb der ausgelieferte Quelltext gelesen, nicht die
Dokumentation und nicht allein die deklarierten Abhängigkeiten. Geprüft sind 27 von 203
verwendeten Paketen, die 87,9\,\% der npm-nutzenden Mandanten erreichen.

\subsubsection{Ergebnis}

713 Mandanten (41,5\,\%) benötigen eine Schnittstelle, die eine sprachbasierte Laufzeit ohne
vermitteltes Binding nicht bereitstellt. Diese Gruppe wird im Folgenden als Hard Floor
bezeichnet. Werden die Mandanten, deren Pakete vollständig außerhalb der geprüften Menge liegen,
pessimistisch als nicht lauffähig gewertet, ergibt sich ein Viabilitätsintervall von 39,0\,\% bis
42,9\,\%.

\begin{figure}[H]
  \centering
  \includegraphics[width=\textwidth]{hard_floor_einordnung.png}
  \caption{Anteil des Hard Floor am Gesamtbestand und Größenverhältnis der beiden
    blockierenden Paketgruppen}
  \label{fig:hard-floor-einordnung}
\end{figure}

Die blockierenden Pakete zerfallen in zwei Gruppen mit deutlich verschiedenem Bedarf, deren
Größen \autoref{tab:blockierende-familien} gegenüberstellt.

\begin{table}[H]
  \caption{Blockierende Paketgruppen im Hard-Floor-Bestand}
  \label{tab:blockierende-familien}
  \small
  \begin{tabularx}{\textwidth}{@{}>{\RaggedRight\arraybackslash}X r r >{\RaggedRight\arraybackslash}p{0.30\textwidth}@{}}
    \hline
    Gruppe & Pakete & Mandanten & stärkster Vertreter \\
    \hline
    Benötigen \ac{HTTP} direkt und umgehen ein bereitgestelltes \texttt{fetch} & 14 & 629 &
      \texttt{axios} mit 579 Mandanten \\
    Benötigen einen rohen \ac{TCP}- oder \ac{TLS}-Socket & 8 & 63 &
      \texttt{mssql} mit 46 Mandanten \\
    \hline
    \multicolumn{2}{@{}l}{Schnittmenge beider Gruppen} & 44 & \\
    \multicolumn{2}{@{}l}{Vereinigung} & 648 & \\
    \multicolumn{2}{@{}l}{Hard Floor gesamt} & 713 & \\
    \hline
  \end{tabularx}
\end{table}

648 der 713 blockierten Mandanten, also 90,9\,\%, scheitern an einem Netzwerk-Binding und an
nichts anderem. Dateisystemzugriff, Prozesserzeugung und dynamische Codeerzeugung tragen nur den
Rest bei, wie \autoref{fig:hard-floor-zusammensetzung} nach blockierender Ursache aufschlüsselt.
Die Kompatibilitätsschranke ist damit eng umrissen, und zwischen den beiden Gruppen liegt ein
Größenunterschied von etwa zehn zu eins.

% PRUEFEN: fig:hard-floor-zusammensetzung ueberschneidet sich inhaltlich mit
% tab:blockierende-familien. Eine von beiden Darstellungen genuegt moeglicherweise.

\begin{figure}[H]
  \centering
  \includegraphics[width=\textwidth]{hard_floor_zusammensetzung.png}
  \caption{Zusammensetzung des Hard Floor nach blockierender Ursache}
  \label{fig:hard-floor-zusammensetzung}
\end{figure}

Die kleinere Gruppe stellt die schwierigere Anforderung. Ihre Vertreter sind überwiegend
Datenbanktreiber mit eigenem binären Protokoll. Bei \texttt{mssql} und dem darunterliegenden
\texttt{tedious} wird die Verbindung im Klartext aufgebaut und erst danach auf \ac{TLS}
umgestellt. Ein bereitgestelltes \texttt{fetch} hilft dieser Gruppe nicht.

Ein zweiter Befund derselben Auswertung betrifft nicht die Lauffähigkeit, sondern die Größe der
gebündelten Artefakte. Von den 14.575 Skriptversionen binden 2.373 überhaupt npm-Pakete ein, und
für diese liegt das fertige Bundle im Median bei 69,3\,KB, im 95.\ Perzentil bei 429,0\,KB und im
Maximum bei 1,41\,MB. Betroffen sind 1.023 der 1.718 Mandanten, also knapp 60\,\%. Gegenüber dem
Quelltext, der meist wenige Kilobyte misst, entspricht das typischerweise einer Aufblähung um den
Faktor zehn bis dreißig, die fast vollständig auf die Abhängigkeiten entfällt. Je Mandant liegt
die Zahl gebündelter Versionen im Median bei zwei und im 95.\ Perzentil bei dreizehn, im
Einzelfall bei 151.

Für die Kompatibilitätsfrage ist diese Größe ohne Belang. Für den Entwurf ist sie es nicht, weil
das Bundle nach \autoref{sec:zustellung} bei jeder Ausführung zugestellt und ausgewertet werden
muss. Sobald die Laufzeit selbst in wenigen Millisekunden bereitsteht, bestimmen diese Kilobyte
den Anfragepfad.

\subsubsection{Grenzen der Auswertung}

Der Datensatz deckt allein die heutige Webhook-Verarbeitung ab. Für das Agentenprofil existiert
im Produkt nichts, was migriert werden müsste, weshalb sich aus diesen Daten dazu keine Aussage
ableiten lässt.

Aufrufzahlen enthält der Datensatz nicht. Ein Mandant mit einem Skript pro Monat wiegt so schwer
wie einer mit hoher Aufruffrequenz, weshalb die Mandantenquote eine Obergrenze der praktischen
Betroffenheit ist und keine Schätzung. Die transitive Hülle ist zudem unvollständig, was das
angegebene Intervall abbildet. Die Auswertung erfolgte aggregiert, ohne Codeausschnitte und ohne
Mandantenkennungen.

\subsection{Angreifermodell und relevante Angriffswege}
\label{sec:angreifermodell}

Das Angreifermodell wird über Fähigkeiten bestimmt und nicht über Identität. Die maßgebliche
Fähigkeit besteht darin, beliebigen JavaScript-Code innerhalb der Ausführungsumgebung eines
Mandanten auszuführen. Sie ist durch die Bauart des Dienstes gegeben und muss nicht erst erlangt
werden. Die Codeausführung ist damit kein Angriffsziel, sondern der Ausgangspunkt, und die
maßgebliche Frage lautet, was der Code von dort aus erreichen kann. Dieselbe Ausgangslage
kennzeichnet die Bewertung von Ausbruchsfähigkeiten aktueller Sprachmodelle, bei der der Zugang
innerhalb des Containers vorausgesetzt und allein der Ausbruch untersucht wird
\autocite{marchand_quantifying_2026}.

Wer den Code verfasst, ist für diese Fähigkeit unerheblich. Nach V1 kann es ein bösartiger Mandant
sein, nach V2 im Agentenprofil auch ein Dritter, dessen Inhalte in den Kontext des Modells
gelangen. Beide Urheber führen auf dieselbe technische Lage. Die Ziele des Angreifers entsprechen
den negierten Schutzzielen aus \autoref{sec:systemkontext}. Alles, was V3 der Plattform zurechnet,
gehört nicht zu seinen Fähigkeiten, also weder physischer Zugriff noch ein Zugang auf dem Host
außerhalb der Ausführungsumgebung noch eine Netzwerkposition im Plattforminneren.

Die verbleibenden Angriffswege lassen sich nach der Schicht ordnen, auf der die ausgenutzte
Ressource geteilt wird, entsprechend der Einordnung von Angriffswegen serverloser Plattformen
nach Anwendungs- und Infrastrukturebene \autocite{marin_serverless_2022}.

\subsubsection{A1, Anwendungsebene: Zustand zwischen Ausführungen}

Weil der Code in denselben Prozess geladen wird, überdauern über \texttt{setTimeout} oder
\texttt{setInterval} registrierte Callbacks das Ende ihrer Ausführung und können Daten späterer
Ausführungen mitlesen \autocite{dvelop_ag_dvelop_2026}. Ebenso überdauern globale Variablen und
Reste im Speicher. Das Zurücksetzen des Dateisystems greift nicht, weil der übertragene Zustand
im Speicher liegt, und das Verwerfen des Modul-Zwischenspeichers greift nur für Zustand auf
Modulebene. Globale Variablen, veränderte Prototypen, registrierte Callbacks und offene
Verbindungen überdauern weiterhin. Da eine Instanz genau einen Mandanten bedient, verläuft dieser Weg innerhalb
eines Mandanten, und seine Bewertung folgt der in \autoref{sec:systemkontext} beschriebenen
Verschiebung der Vertrauenseinheit.

Derselbe geteilte Prozess trägt einen zweiten Weg. Eine Ausführung kann Ressourcen so erschöpfen,
dass die Instanz für folgende Ausführungen unbrauchbar wird. Der Ausgangszustand belegt das an
einem Vorfall, bei dem ein Verfügbarkeitstest im Minutentakt eine Instanz dauerhaft aktiv hielt
und ein Speicherleck ihren Verbrauch bis zur Unbrauchbarkeit steigen ließ.

\subsubsection{A2, Sprachebene: geteilte Laufzeit im selben Prozess}

% GRUNDLAGEN: Mikroarchitekturelle Seitenkanaele und Spectre gehoeren als Begriff
% nach Kapitel 2. Hier bleibt nur die Folgerung fuer das Bedrohungsmodell.

Der zweite Weg greift eine Isolationsgrenze an, die innerhalb eines Prozesses verläuft und
Angriffe unterhalb der Softwareebene nicht abfangen kann. Über einen entfernten Timer ließ sich
aus einem ko-lokierten Cloudflare Worker ein JSON Web Token mit bis zu zwölf Bit pro Sekunde bei
einer Genauigkeit von rund 99 Prozent auslesen, und zwar auf einer produktiven Plattform, die V8
Isolates einsetzt \autocite{schwarzl_remote-timer-as--service_2026}. Denselben Befund hatte zuvor
die Einführung der Site Isolation in Chrome, wo eine sprachbasierte Trennung zweier Origins im
gemeinsamen Prozess nach Bekanntwerden von Spectre nicht mehr trug und durch eine Prozessgrenze
ersetzt wurde \autocite{reis_site_2019}. Cloudflare benennt die daraus folgenden Restrisiken für
V8 Isolates ausdrücklich \autocite{cloudflare_security_2026}. Eine sprachbasierte Grenze allein
bildet damit keine hinreichende Mandantengrenze.

\subsubsection{A3, Infrastrukturebene: Sandbox-Laufzeit und Host-Kernel}

Der dritte Weg zielt auf den Ausbruch aus der Ausführungsumgebung. Generische Containerisierung
trägt diese Grenze nur eingeschränkt. Das \ac{NIST} führt den geteilten Kernel als eigene
Angriffsfläche und hält fest, dass ein Container eine Sandbox nicht vollwertig ersetzt
\autocite{souppaya_application_2017}. Eine Untersuchung an neun öffentlichen Ausbruchs-Exploits
zeigt, dass eine Einschränkung der Containerrechte über User Namespaces sieben davon abwehrt,
zwei aber erfolgreich bleiben \autocite{reeves_towards_2021}. Die Fallstudie begegnet dem bereits
über gVisor, dessen Sentry die Systemaufrufe des Gastcodes selbst beantwortet und die zugängliche
Oberfläche des Host-Kernels verkleinert \autocite{the_gvisor_authors_security_2026}. Die
Angriffsfläche verschwindet dadurch nicht, sie verlagert sich, denn der Sentry ist selbst in
Software implementiert und angreifbar, was eine gezielte Fuzzing-Untersuchung belegt
\autocite{li_g-fuzz_2024}.

Der Weg bleibt damit relevant, wenn auch auf höherem Niveau. Ein Agent, der innerhalb eines
Containers über Codeausführung verfügt, kann Fehlkonfigurationen, Rechtevergabefehler,
Kernel-Schwachstellen sowie Schwächen der Laufzeit und der Orchestrierung selbst auffinden und
ausnutzen \autocite{marchand_quantifying_2026}. Diese Untersuchung bezieht sich auf Container ohne
vorgeschaltete Sandbox-Laufzeit und ist deshalb nicht unmittelbar auf die Fallstudie übertragbar.
Sie beschreibt jedoch den Angreifertyp, dem die Grenze standhalten muss, sobald Codeausführung im
Container zum bestimmungsgemäßen Betrieb gehört. Aus einem menschlichen Angreifer mit
vorbereitetem Exploit wird der ausgeführte Workload selbst.

\subsection{Kriterienkatalog}
\label{sec:kriterienkatalog}

Der Katalog übersetzt die Ergebnisse der vorangegangenen Abschnitte in einen Maßstab. Jedes
Kriterium ist auf seine Herkunft zurückführbar, also auf einen Angriffsweg, eine
Vertrauensannahme, die Bestandsanalyse oder eine Betriebsanforderung. Kriterien ohne solche
Herkunft werden nicht aufgenommen.

Unterschieden werden zwei Arten. \emph{Muss-Kriterien} beschreiben Eigenschaften, ohne die ein
Ausführungsmodell ausscheidet. Sie werden einzeln geprüft und nicht verrechnet, und sie
entscheiden die Vorauswahl in \autoref{sec:vorauswahl}. \emph{Soll-Kriterien} beschreiben
Eigenschaften, in denen sich Modelle graduell unterscheiden. Sie scheiden nichts aus, sondern
bilden den gewichteten Maßstab, an dem \autoref{sec:Evaluation} das fertige Artefakt gegen den
Ausgangszustand bewertet. Diese Trennung ist Voraussetzung dafür, dass das Artefakt scheitern
kann, denn ein Modell mit gutem Nutzwert, das ein Muss-Kriterium verletzt, ist keine Lösung.

\subsubsection{Muss-Kriterien}

\begin{table}[H]
  \caption{Muss-Kriterien, Herkunft und Nachweis}
  \label{tab:musskriterien}
  \small
  \begin{tabularx}{\textwidth}{@{}l >{\RaggedRight\arraybackslash}X >{\RaggedRight\arraybackslash}p{0.18\textwidth} >{\RaggedRight\arraybackslash}p{0.18\textwidth}@{}}
    \hline
    Nr. & Kriterium & Herkunft & Nachweis \\
    \hline
    M1 & \textbf{Prozessgrenze je Mandant.} Zwischen zwei Mandanten verläuft eine Grenze auf
         Prozess- oder Virtualisierungsebene. Eine sprachbasierte Grenze allein genügt nicht. &
         A2 & Architektur\-prüfung \\[2pt]
    M2 & \textbf{Zustandsfreiheit je Ausführung.} Kein Zustand einer Ausführung ist einer
         späteren zugänglich, auch nicht innerhalb eines Mandanten. Umfasst Prozessspeicher,
         globale Variablen, registrierte Callbacks, offene Verbindungen und Dateisystem. &
         A1, Verschiebung der Vertrauens\-einheit & Prüffälle T1 bis T6 \\[2pt]
    M3 & \textbf{Keine größere Kernel-Oberfläche.} Die dem Gastcode zugängliche Oberfläche des
         Host-Kernels wächst gegenüber dem Ausgangszustand nicht. &
         A3 & Konfigurations\-vergleich \\[2pt]
    M4 & \textbf{Vermittelter Zugriff.} Netz und Dateisystem ausschließlich über Bindings, die
         die Ausführungsumgebung vermittelt. Keine Beschränkung beruht auf Mitwirkung des
         Modells. & V5 & Prüffälle T7 und T8 \\[2pt]
    M5 & \textbf{Verfügbarkeit bei Knotenupdates.} Die Zeit ohne Ausführungsfähigkeit während
         eines rollierenden Knotenupdates überschreitet den festgelegten Wert nicht. &
         Betriebs\-anforderung, vier Updates je Jahr & Messung, Zielwert festzulegen \\
    \hline
  \end{tabularx}
\end{table}

M3 ist gegen den Ausgangszustand als Baseline formuliert, weil dieser mit gVisor bereits einen
User-Space-Kernel einsetzt. Ein Modell, das diese Schicht aufgibt, um Startzeit zu gewinnen, wäre
auf der Infrastrukturebene ein Rückschritt, auch bei besserem Nutzwert. M4 folgt unmittelbar aus
V5, denn die Schnittstelle zur Ausführungsumgebung ist der einzige Durchsetzungspunkt, den die
Plattform kontrolliert.

\subsubsection{Soll-Kriterien und Gewichtung}

\begin{table}[H]
  \caption{Soll-Kriterien, Messgrößen und Gewichte, absteigend nach Gewicht}
  \label{tab:sollkriterien}
  \small
  \begin{tabularx}{\textwidth}{@{}l >{\RaggedRight\arraybackslash}p{0.23\textwidth} >{\RaggedRight\arraybackslash}X r@{}}
    \hline
    Nr. & Kriterium & Messgröße & Gewicht \\
    \hline
    S1 & Kaltstartlatenz & Zeit bis zum Ausführungsbeginn ohne warme Instanz, P95 & 0,30 \\
    S6 & Kompatibilität des Bestands & Anteil unverändert lauffähiger Mandanten und
         Skripte & 0,25 \\
    S2 & Warme Latenz & Ausführungslatenz bei warmer Instanz, P50 bis P99 & 0,15 \\
    S3 & Speicherbedarf & Bedarf je Ausführung und je Mandant im Ruhezustand & 0,10 \\
    S7 & Aufwand für den Rest & Aufwand, den nicht lauffähigen Anteil zu überführen & 0,10 \\
    S4 & Durchsatz & Ausführungen je Sekunde bei gesättigter Last & 0,05 \\
    S5 & Integrierbarkeit & Anteil der an vorhandene Bausteine abgebbaren Aufgaben & 0,05 \\
    \hline
    \multicolumn{3}{@{}l}{Summe} & 1,00 \\
    \hline
  \end{tabularx}
\end{table}

Die Gewichte sind eine begründete Setzung und kein Messergebnis. S1 erhält das höchste Gewicht,
weil die Kaltstartzeit im Agentenprofil in die Wartezeit der anfragenden Person eingeht und sich
über die Aufrufe einer Aufgabe summiert. S6 folgt darauf, weil \autoref{sec:bestandsanalyse} die
Kompatibilität als stärkste Gegenkraft zu einem Wechsel des Ausführungsmodells ausweist. Die
übrigen Kriterien wiegen leichter, weil sie entweder die Ausführung selbst betreffen, ohne die
Wartezeit im gleichen Maß zu prägen wie S1, oder weil sie Restgrößen abdecken, deren Größenordnung
bereits durch S6 begrenzt ist. Die Bewertung je Kriterium erfolgt auf einer Skala von null bis
fünf, und der Nutzwert ist die mit den Gewichten gebildete Summe. Die Zuordnung von Messwerten zu
Punktwerten legt \autoref{sec:Evaluation} offen, damit die Bewertung nicht von der Darstellung
abhängt.

\subsubsection{Prüffallkatalog}

M2 und M4 lassen sich nicht durch Architekturbetrachtung allein belegen. Für sie wird ein Katalog
von Prüffällen aufgestellt, der aus den Angriffswegen A1 bis A3 abgeleitet ist. Jeder Prüffall
wird gegen den Ausgangszustand und gegen das Artefakt ausgeführt, sodass die Gegenüberstellung
zugleich belegt, welche Eigenschaften der Ausgangszustand bereits erfüllt.

\begin{table}[H]
  \caption{Prüffälle zur Zustandsfreiheit und zur Durchsetzung der Fähigkeitsbeschränkung}
  \label{tab:prueffaelle}
  \small
  \begin{tabularx}{\textwidth}{@{}l >{\RaggedRight\arraybackslash}X l l@{}}
    \hline
    Nr. & Prüffall & Kriterium & Quelle \\
    \hline
    T1 & Ein über \texttt{setTimeout} registrierter Callback überdauert das Ende seiner
         Ausführung und greift auf Daten einer späteren Ausführung zu. & M2 & beobachtet \\
    T2 & Eine in Ausführung $n$ gesetzte globale Variable ist in Ausführung $n+1$ lesbar. &
         M2 & abgeleitet \\
    T3 & Eine Veränderung eines eingebauten Prototyps oder des Module Cache wirkt in
         Ausführung $n+1$ fort. & M2 & abgeleitet \\
    T4 & Eine in Ausführung $n$ abgelegte Datei ist in Ausführung $n+1$ lesbar. & M2 &
         abgeleitet \\
    T5 & Eine in Ausführung $n$ geöffnete Netzwerkverbindung ist in Ausführung $n+1$ nutzbar.
         & M2 & abgeleitet \\
    T6 & Eine Ausführung erschöpft den Speicher und macht die Instanz für folgende
         Ausführungen unbrauchbar. & M2 & beobachtet \\
    T7 & Ausgehender Netzwerkverkehr wird an dem vermittelten Binding vorbei aufgebaut. &
         M4 & abgeleitet \\
    T8 & Auf das Dateisystem wird an dem vermittelten Binding vorbei zugegriffen. & M4 &
         abgeleitet \\
    \hline
  \end{tabularx}
\end{table}

T1 und T6 bilden Vorgänge nach, die im Ausgangszustand dokumentiert sind, und sind damit keine
konstruierten Laborfälle. Für T4 gilt das nicht. Der Ausgangszustand setzt das temporäre
Verzeichnis vor jeder Ausführung zurück, weshalb eine Datei aus Ausführung $n$ in Ausführung
$n+1$ gerade nicht mehr vorhanden ist. T4 prüft also, ob dieser Reset auch im Artefakt greift.
Die übrigen Prüffälle ergänzen den Katalog um Wege, die aus derselben geteilten Ressource folgen,
aber noch nicht beobachtet wurden.

Damit steht fest, wodurch die Unterfragen aus \autoref{sec:forschungsfragen} beantwortet werden.
U3 prüft M1 bis M4 über Prüffallkatalog und Architekturprüfung, U4 erhebt die Messgrößen S1 bis
S4 sowie M5, und die Muss-Kriterien entscheiden in \autoref{sec:Entwurf} über die Technologie.
